This paper set out to measure the menu of Python acceleration techniques under conditions that
let them be compared: eight questions, one measurement protocol, one interpreter provenance,
one documented host, and code for every number. What comes out of it is a table indexed
by workload rather than by tool, and the sections below are the entries that a reader should
take away from it.

\textbf{What the technology buys.} Compiled code --- by any of the routes tested --- is one
to two orders of magnitude faster than the CPython interpreter on compute-bound kernels
($31$--$\meas{ratio:b1_compute-314:arraysum/cpython_loop|b1_compute-314:arraysum/codon_pyext_ptr}{32}\times$ on an array sum, $52$--$\meas{ratio:b1_compute-314:mandelbrot/cpython_loop|b1_compute-314:mandelbrot/cython}{62}\times$ on a mandelbrot grid, $46$--$\meas{ratio:b1_compute-314:matmul/cpython_loop|b1_compute-314:matmul/numba}{128}\times$ on a
naive matrix product). Which route you take matters far less than taking one: on the
scalar reduction they agree to within $\meas{pct:b1_compute-314:arraysum/codon_jit|b1_compute-314:arraysum/codon_pyext_ptr}{4.4}\%$, and on the other two kernels they differ by
$\meas{pct:b1_compute-314:mandelbrot/codon_jit|b1_compute-314:mandelbrot/cython}{21}\%$ and by $\meas{ratio:b1_compute-314:matmul/codon_jit|b1_compute-314:matmul/numba}{2.8}\times$, in
both directions and without a consistent winner. Against a factor of thirty to a hundred, speed
does not decide between Cython, Codon, Numba, C and Rust; what does --- build complexity,
debuggability, code ownership --- this study does not measure. It is a real spread, not
a tie, and a project whose whole cost is one kernel should measure its own. Removing the GIL delivers
$\meas{ratio:b4_threads-ft:py_arith/ft/T1_control|b4_threads-ft:py_arith/ft/T8}{5.60}\times$ on arithmetic and $\meas{ratio:b4_threads-ft:py_branchy/ft/T1|b4_threads-ft:py_branchy/ft/T8}{6.05}\times$ on branch-heavy code across eight threads,
and a \code{PYTHON\_GIL=1} control on the same binary --- which instead loses $14$--$\meas{saving:b4_threads-ft_gil:py_arith/ft_gil/T1_control|b4_threads-ft_gil:py_arith/ft_gil/T8}{20}\%$ at
eight threads --- attributes \emph{that scaling} to the lock rather than the build; the price of
obtaining it runs the other way, and is taken up below.

\textbf{Where the mechanism at work was not the one in the technique's name.} The largest levers in the compute
study are the decision to compile at all and, where the operation already exists in a library,
a vendor BLAS; the flags underneath those decisions
behave less uniformly than their names suggest. Disassembling every flag variant settles what
each one actually changes: \code{-O2} and \code{-O3} emit byte-identical code on two of the
three kernels and differ by a single instruction on the third, with run times to match, so the
step between them buys nothing here. Only \code{-O3 -march=native} emits different code
everywhere, and it is the flag whose effect runs in both directions (\S\ref{sec:rq1b}).
Turning the vectoriser off under that flag changed neither the emitted code nor
the time on any kernel, and SIMD instructions appear only once fast-math is granted, so the
architecture flag is not buying vectorisation here. The floating-point contract is worth
$\meas{ratio:b1_flags-314:arraysum/O3native|b1_flags-314:arraysum/O3native_ffast}{2.42}\times$ on the reduction and is not really a flag either: four independent accumulators
written into the C source, at plain \code{-O2} with no fast-math anywhere, are $\meas{pct:b1_flags-314:arraysum/O3native_ffast|b1_flags-314:arraysum/accum_acc4}{13}\%$ faster than
\code{-ffast-math} itself, and eight accumulators are worse than four. The row that would read
``use native flags'' reads instead ``decide what the arithmetic is allowed to do, know that you
are deciding it, and disassemble the result on the target you ship''. On irregular code,
branching is not what defeats a compiler (Cython and C are within half a per cent on our graph
traversal, and the native routes lead by at most $\meas{pct:b2_branchy-314:tokenize/codon_jit|b2_branchy-314:tokenize/c_ctypes}{17}\%$ on a byte scanner) ---
\emph{allocation} is, and that is where C and Rust pull ahead by $3.7$--$\meas{ratio:b2_branchy-314:binarytrees/numba|b2_branchy-314:binarytrees/rust_ctypes}{5.3}\times$.

\textbf{Where the answer is a distribution rather than a number.} Interpreter progress from
3.10 to 3.14 is real and uneven: 3.11 delivered the step (attribute reads $\meas{ratio:b3_runtime-v310:attr_get/v310|b3_runtime-v311:attr_get/v311}{1.84}\times$, method
calls $\meas{ratio:b3_runtime-v310:method_call/v310|b3_runtime-v311:method_call/v311}{1.76}\times$ against 3.10), 3.12 gave part of it back --- it is slower than 3.11 on 14 of
our 16 operations, all 14 significant, although it is the release in which \code{None} became
immortal --- 3.13 changed little, and 3.14 recovered to $\meas{ratio:b3_runtime-v310:attr_get/v310|b3_runtime-v314:attr_get/v314}{1.97}\times$ and $\meas{ratio:b3_runtime-v310:method_call/v310|b3_runtime-v314:method_call/v314}{2.04}\times$, while
building a dict literal on 3.14 runs at $\meas{ratio:b3_runtime-v310:create_dict/v310|b3_runtime-v314:create_dict/v314}{0.74}\times$ its 3.10 rate. At the build level the
useful entry is the combination rather than any single option: PGO on its own is worth
$\meas{geomean:b5_buildcfg-plain|b5_buildcfg-pgo}{1.02}\times$, PGO with \emph{full} LTO $\meas{geomean:b5_buildcfg-plain|b5_buildcfg-pgo_ltofull}{1.10}\times$, and adding \code{-march=native} to that
$\meas{geomean:b5_buildcfg-plain|b5_buildcfg-pgo_ltofull_native}{1.12}\times$; attribute writes are the one operation that loses in all six configurations, at
$0.88$--$\meas{ratio:b5_buildcfg-plain:attr_set/plain|b5_buildcfg-pgo:attr_set/pgo}{0.98}\times$. CinderX resists a one-word verdict in the same way. It builds in both of
its configurations with zero compiler errors, its JIT force-compiled all eight functions we
handed it without a failure, and every kernel passed the correctness gate. Its runtime then
costs $1$--$\meas{pct:b6_cinderx-fork_cinderx:mandelbrot/fork_cinderx|b6_cinderx-stock:mandelbrot/stock}{33}\%$ on our six untyped kernels in either configuration, which puts that cost in
the code path rather than in a build flag, and its JIT repays it on all six at
$0.33$--$\meas{ratio:b6_cinderx-fork_jit:binarytrees/fork_jit|b6_cinderx-stock:binarytrees/stock}{0.79}\times$ stock --- but only when compilation follows specialisation rather than
preceding it. Its lightweight frames moved our probes by under $\meas{pct:b6_cinderx-stock:traceback_depth_50/stock|b6_cinderx-fork:traceback_depth_50/fork}{3}\%$ on the bare
fork. Its parallel collector moved its probe the wrong way while the worker's long-lived heap
was visible to it ($\meas{ratio:b6_gcscale-fork_cinderx:gc_collect/fork_cinderx/state_visible_par6|b6_gcscale-fork_cinderx:gc_collect/fork_cinderx/state_visible_serial}{1.48}\times$ slower) and the right way only once \code{immortalize\_heap()}
had taken that heap out of view ($\meas{ratio:b6_gcscale-fork_cinderx:gc_collect/fork_cinderx/immortal_serial|b6_gcscale-fork_cinderx:gc_collect/fork_cinderx/immortal_par6}{1.06}\times$ faster), which makes heap composition the variable
rather than the thread count. Its one large win (Static Python, $\meas{ratio:b6_static-fork_static_adaptive:matmul96_int/fork_static_adaptive/boxed_jit|b6_static-fork_static_adaptive:matmul96_int/fork_static_adaptive/static_jit}{11.4}\times$ on top of the JIT) is reachable only through
a boot sequence documented only in a tutorial and only for rewritten typed code. Measuring the adaptive
configuration was itself work: the specialising path for static opcodes had been carried to 3.14
without the handlers it specialises into, so Static Python raised on its first specialisation
until we supplied them --- and once supplied, the flag takes $\meas{saving:b6_static-fork_static_adaptive:matmul96_int/fork_static_adaptive/static_interp|b6_static-fork_static:matmul96_int/fork_static/static_interp}{43}\%$ off interpreted Static
Python and nothing at all off the compiled kind, which is the clearest statement in the study
that primitive opcodes are an intermediate representation for a JIT rather than a faster way to
interpret.

\textbf{The costs, without which a multiplier is not yet a decision.} The numbers that decide a
choice and are seldom stated alongside it: the point at which a JIT pays for itself, which moves from 7 calls to 7335
depending only on input size, against one observed $\meas{fact:b1_breakeven-314:facts.numba_compile_s:ms}{215}$\,ms of compile latency; the cost of
an FFI boundary, $59$--$\meas{med:b1_compute-314:call_overhead/c_ctypes:ns}{166}$\,ns per crossing (\S\ref{sec:rq2}); the single-thread price of free-threading,
$1.05$--$\meas{ratio:b4_threads-ft:single_thread_arith/ft|b4_threads-gil:single_thread_arith/gil}{1.14}\times$, which the free-threaded \emph{build} charges whether or not the lock is
off and which the second thread repays; and the cost of Static Python, which is
$\meas{ratio:b6_static-fork_static:matmul96_int/fork_static/static_interp|b6_static-fork_static:matmul96_int/fork_static/boxed_python}{10.2}\times$ \emph{slower} than boxed Python if its JIT is not also enabled, and still
$\meas{ratio:b6_static-fork_static_adaptive:matmul96_int/fork_static_adaptive/static_interp|b6_static-fork_static_adaptive:matmul96_int/fork_static_adaptive/boxed_python}{5.8}\times$ slower with the specialising interpreter path switched on. Two further results
belong here because they are cost-dominated. PyPy returns $3.2$--$\meas{ratio:b1_compute-314:matmul/cpython_loop|b1_compute-pypy311:matmul/cpython_loop}{42.8}\times$ on unmodified
pure-Python code --- the largest factor here that asks for no edit. But the one kernel whose
\code{ctypes} call hands it a two-megabyte buffer costs $\meas{ratio:b2_branchy-pypy311:tokenize/c_ctypes|b2_branchy-314:tokenize/c_ctypes}{5.3}\times$ more there than on CPython
($118.5$ against $\meas{med:b2_branchy-314:tokenize/c_ctypes:ms}{22.3}$\,ms), while the calls that pass only integers are within $\meas{pct:b2_branchy-314:binarytrees/rust_ctypes|b2_branchy-pypy311:binarytrees/rust_ctypes}{7}\%$:
``switch runtime'' and ``move the hot loop to C'' can exclude each other. The second is CPython's
own experimental JIT, a $\meas{geomean:b5_buildcfg-plain|b5_buildcfg-jit}{1.04}\times$ proposition overall that is not free either. It improves 10
of our 16 operations and slows 6, integer arithmetic and attribute writes to $0.88$--$\meas{ratio:b5_buildcfg-plain:int_arith/plain|b5_buildcfg-jit:int_arith/jit}{0.89}\times$.

\textbf{On method.} Three of our own results were wrong before the protocol caught them
(\S\ref{sec:method}), which is the whole argument for having one. What caught them was cheap:
delegating the timing to \code{pyperf}, alternating measurement where a claim turns on a few per
cent, running the correctness gate in the process that will do the timing, and probing the host
rather than the program.

\textbf{Practically}, composition survives as a principle with a sharper rule attached: on our
mixed pipeline, Numba on a stage worth \meas{share:b7_pipeline-gil:stage_aggregate/gil/v0_pure|b7_pipeline-gil:end_to_end/gil/v0_pure}{3.6}\% of the run is worth $\meas{ratio:b7_pipeline-gil:end_to_end/gil/v0_pure|b7_pipeline-gil:end_to_end/gil/v1_numba}{1.01}\times$; moving the
classification stage --- half the run time --- into C is worth $\meas{ratio:b7_pipeline-gil:end_to_end/gil/v0_pure|b7_pipeline-gil:end_to_end/gil/v2_numba_native}{1.35}\times$; and spreading the work over
four threads on top of that gives $\meas{ratio:b7_pipeline-gil:end_to_end/gil/v0_pure|b7_pipeline-gil:end_to_end/gil/v3_native_threads}{0.72}\times$ under the GIL and $1.26$--$\meas{ratio:b7_pipeline-ft:end_to_end/ft/v0_pure|b7_pipeline-ft:end_to_end/ft/v3_native_threads}{1.28}\times$ on the
free-threaded build, still below the $\meas{ratio:b7_pipeline-gil:end_to_end/gil/v0_pure|b7_pipeline-gil:end_to_end/gil/v2_numba_native}{1.35}\times$ the one well-aimed change already
delivers. Composition is not additive; it is profile-directed, and stacking past the profile
takes performance back. Optimisation decisions should be indexed by what the workload
does --- loop over arrays, allocate objects, traverse a graph, start up, spread across cores ---
and not by the name of the tool, because within a workload class the spread between tools
was small next to the decision to use one at all, while across classes the same tool differed by
an order of magnitude.
Table~\ref{tab:decision} collects those entries, and every number in it is reproducible from the
accompanying artifact with two commands.
